% ==============================================================
%  Chapter 4: AI for Cloud Engineers
% ==============================================================
\chapter{AI for Cloud Engineers}
\label{ch:cloud}

Cloud engineering sits at the intersection of architecture, automation, cost management,
and security.
AI helps most in the early design phase (generating options and tradeoffs), in
translating requirements into policy and configuration, and in diagnosing unfamiliar
services.
It cannot know your account structure, existing dependencies, or budget constraints
unless you provide them.

\section{Architecture Design and Review}
\label{sec:arch-design}

AI is an excellent architecture sounding board.
It knows AWS, Azure, and GCP well across compute, storage, networking, and managed
services.
Use it to generate options, stress-test a design, and surface considerations you missed.

\textbf{High-value architecture prompts:}
\begin{itemize}
  \item \textit{``I need to process 50,000 events per second with at most 5 seconds end-to-end latency. My team knows Python. What are three architectural approaches on AWS, with tradeoffs for cost, operational complexity, and scalability?''}
  \item \textit{``Review this architecture diagram description for single points of failure, over-provisioned services, and missing observability components.''}
  \item \textit{``What does the AWS Well-Architected Framework say about data transfer costs between services in different AZs?''}
\end{itemize}

\begin{tipbox}
Give AI your constraints explicitly: team skills, existing services, budget ceiling,
RTO/RPO requirements, and compliance framework.
Without constraints, AI generates architectures that are technically correct but
organisationally impractical.
\end{tipbox}

\section{Cost Optimisation}
\label{sec:cost}

Cost optimisation involves pattern recognition across pricing models, usage data, and
service alternatives --- a natural AI use case.

\textbf{Effective cost-review prompts:}
\begin{itemize}
  \item \textit{``My EC2 instances run at 15\% average CPU. What rightsizing, Reserved Instance, or Savings Plan options should I consider?''}
  \item \textit{``We spend \$8,000 per month on data transfer. What architectural changes typically reduce cross-region egress costs?''}
  \item \textit{``Compare Lambda vs Fargate vs ECS on EC2 for a workload that runs 10,000 short (30 s) tasks per day, with cold start tolerance.''}
\end{itemize}

\begin{warnbox}
AI uses pricing data from its training set, which may be out of date.
Always verify current pricing on the cloud provider's pricing calculator before
presenting cost estimates to stakeholders.
\end{warnbox}

\section{Security and IAM}
\label{sec:cloud-security}

Cloud IAM is one of the most error-prone areas in cloud engineering --- and one of the
most valuable AI use cases.
AI can audit policies, generate least-privilege alternatives, and translate compliance
requirements into concrete controls.

\textbf{IAM audit workflow:}
\begin{enumerate}
  \item Export the IAM policy JSON from your cloud console.
  \item Paste it and ask: \textit{``Identify any permissions violating the principle of
        least privilege, and suggest a tighter replacement policy.''}
  \item Ask: \textit{``What is the blast radius if this role is compromised?''}
  \item Ask: \textit{``Does this policy expose any data exfiltration risk?''}
\end{enumerate}

\textbf{Generating service control policies and guardrails:}
\begin{itemize}
  \item \textit{``Write an AWS SCP that prevents any IAM user from disabling CloudTrail, deleting S3 buckets with versioning, or creating resources outside eu-west-1.''}
  \item \textit{``Generate an Azure Policy that enforces required tags on all resources.''}
\end{itemize}

\section{Multi-Cloud and Migration}
\label{sec:multi-cloud}

AI understands the conceptual mapping between cloud providers:
AWS $\leftrightarrow$ Azure $\leftrightarrow$ GCP equivalents for compute, storage,
and managed services.

\textbf{Migration prompts:}
\begin{itemize}
  \item \textit{``Map these AWS services to Azure equivalents: EC2, RDS PostgreSQL, S3, SQS, Lambda, CloudWatch.''}
  \item \textit{``What are the main operational differences I should expect when migrating from AWS EKS to Azure AKS?''}
  \item \textit{``Write a migration plan outline for moving from self-managed Kubernetes to a managed service. What are the top 5 risks?''}
\end{itemize}

\begin{tipbox}
Use AI to generate the migration task list and risk register as a first draft.
Your cloud architects then validate feasibility against your actual environment.
AI cannot assess undocumented dependencies --- that's your team's value.
\end{tipbox}

\section{Generating Cloud Documentation}
\label{sec:cloud-docs}

Cloud engineers maintain architecture diagrams, runbooks, and disaster recovery
documents.
AI accelerates first-draft creation significantly:

\begin{itemize}
  \item \textit{``Write a runbook for rotating the RDS master password with zero downtime.''}
  \item \textit{``Draft a disaster recovery plan for a three-tier web application hosted on AWS, with RTO 4 hours and RPO 1 hour.''}
  \item \textit{``Generate a service dependency description for this architecture [describe].''}
\end{itemize}
